% ==========================================================================================
% APPENDIX VF (2026 reconstruction): The Master-Scale Chain Cl(6) -> v_F -> C -> M*
% The original VF source was never committed (PDF-only); reconstructed here from the
% compiled report's content, with the 2026 condition note (GAP-N4c) attached.
% ==========================================================================================
\section{Appendix VF: The Master-Scale Chain $Cl(6)\to v_F\to C\to M^*$}
\label{app:VF}

\begin{mdframed}[backgroundcolor=yellow!8, frametitle={Status note (2026)}]
This chain is the \emph{anchor} of the absolute scale used in Section~\ref{sec:genesis-tower}
($M_{\rm stage}=M^*$). It is reconstructed from the February compilation (its source file was
never committed) and is carried as an \textbf{inherited derivation with an open independent
audit} (register item GAP-N4c, Appendix~AC). Nothing else in the Genesis chain depends on its
details --- only the absolute normalization does; mass \emph{ratios} are anchor-independent.
\end{mdframed}

\subsection*{VF.1 Chirality Reduction ($D_e=5$)}
In $Cl(6)$ the chirality operator $\Gamma_7=i\Gamma_1\cdots\Gamma_6$ satisfies
$\Gamma_7^2=1$, and on the chiral subspace
$\Gamma_6=(+i)\,\Gamma_5\Gamma_4\Gamma_3\Gamma_2\Gamma_1$ (verified numerically, error $=0$):
only $D_e=5$ of the six generators are algebraically independent hopping channels.

\subsection*{VF.2 Equipartition ($t=1/5$)}
By $\mathrm{Spin}(6)$ isotropy (Schur) and bandwidth normalization, the tight-binding hopping
is $t=1/D_e=1/5$.

\subsection*{VF.3 Abrikosov locking ($q=6$, $k_F=5/6$)}
The minimum-energy vortex configuration is the triangular lattice with $C_6$ symmetry
(Abrikosov ratio $\beta_A=1.1596<1.1803$ of the square lattice): holonomy $\mathbb{Z}_6$,
flux denominator $q=6$, edge-locking $k_F=1-1/q=5/6$.

\subsection*{VF.4 The exact identity and the DOS constant}
\begin{equation}
v_F \;=\; \frac{2}{D_e}\,\sin\frac{\pi}{q} \;=\; \frac{2}{5}\cdot\frac{1}{2} \;=\; \frac{1}{5}
\qquad(\text{numerically exact}),
\end{equation}
\begin{equation}
C \;=\; g\,\frac{L_F}{4\pi^2}\,\frac{2}{v_F}
\;=\; 4\cdot\frac{5\pi/3}{4\pi^2}\cdot 10 \;=\; \frac{50}{3\pi}\;\approx\;5.305,
\end{equation}
with $g=4$ (Kramers $\times$ particle--hole; Appendix~AB item AB.3) and $L_F=2\pi k_F=5\pi/3$.
An earlier draft additionally quoted a ``lattice cross-check $C=5.339$ ($0.7\%$ finite-size
effect)''; the 2026 audit found no generating code for that number and no simple estimator
artifact that cleanly reproduces it, and it is \textbf{struck} (VF.6, item 2026-08-14b).

\subsection*{VF.5 The master scale}
With $X=3/(2\alpha)\approx205.53$ and $g_{\rm eff}=C/X\approx0.0258$, dimensional
transmutation from the Planck cutoff gives
\begin{equation}
M^{*} \;=\; \Lambda_{\rm UV}\, e^{-1/g_{\rm eff}} \;\approx\; 365.26\ \text{GeV},
\qquad
m_\tau \;=\; \frac{2\alpha}{3}\,M^{*} \;=\; 1776.9\ \text{MeV}
\end{equation}
--- the last relation reproducing the tau mass as a consistency output rather than an input.
The 17-order Planck--electroweak hierarchy is the RG exponential $e^{-1/0.0258}$, the same
mechanism proven exactly at the Layer-0 level (Section~\ref{sec:genesis-condensation}).


\paragraph{Data provenance (2026-08-13 verification).} The input constants were
re-verified against pdgLive during the standardization campaign:
$m_\tau = 1776.93\pm0.09$\,MeV (the repository previously carried the older world average
$1776.86$), $m_W = 80.3692\pm0.0133$\,GeV, $m_Z = 91.1879\pm0.0020$\,GeV. The verified
$m_\tau$ shifts $M^{*}$ from $365.24$ to $365.26$\,GeV ($+0.004\%$) --- inside every quoted
precision; legacy occurrences of $365.24$ elsewhere in this report are within that band.

\subsection*{VF.6 Independent audit results (2026-08-14, GAP-N4c partially executed)}
The lab's independent audit (\texttt{experiments/verification/vf\_chain\_audit.py}, log
\texttt{vf\_chain\_audit\_20260814}) established:
\begin{itemize}
\item \textbf{Verified [THM/NUM]:} the $Cl(6)$ chirality reduction was reconstructed from
scratch (independent $8\times8$ gamma representation): Clifford relations, $\Gamma_7^2=1$,
and $\Gamma_6=+i\,\Gamma_5\Gamma_4\Gamma_3\Gamma_2\Gamma_1$ on the chiral subspace all hold
to machine precision, confirming $D_e=5$. The identities $v_F=1/5$ and $C=50/(3\pi)$ are
exact; the chain reproduces $m_\tau$ to $0.5\%$, and inverting it, the $C$ \emph{required}
by the PDG $\tau$ mass agrees with $50/(3\pi)$ to $0.012\%$.
\item \textbf{Fragility [NUM]:} the transmutation exponent amplifies upstream errors by
$|d\ln M^*/d\ln C| = X/C = 38.7$. Consequently the ``lattice cross-check'' value
$C=5.339$ quoted in VF.4 would shift $M^*$ by $+28\%$ and \emph{miss} the $\tau$ mass
entirely --- only the continuum value works.
\item \textbf{Irreproducible [FAIL]:} no code generating the lattice value $5.339$ exists in
the repository; the archived script \texttt{verify\_C\_band\_structure.py} merely
re-evaluates the same analytic formula (circular). The finite-size claim is therefore
unverifiable as stated.
\item \textbf{Unfixed conventions [OPEN]:} the BCS prefactor ($1$ vs $2$: factor $2.0$ in
$M^*$), the cutoff identification (full vs reduced Planck mass: factor $5.0$), and
$g=4$ are load-bearing choices not fixed by any stated principle.
\end{itemize}
\textbf{Audit verdict:} the arithmetic is sound and the $0.012\%$ inversion agreement is
striking, but with three unfixed discrete choices the chain retains enough freedom that the
agreement cannot yet be claimed as a derivation (Anti-Hardcode Law).

\paragraph{2026-08-14b: the decisive lattice computation (executed).}
The reproducible lattice/DOS computation requested above now exists
(\texttt{experiments/verification/lattice\_C\_computation.py}, log
\texttt{lattice\_C\_computation\_20260814}; contour-integral method validated to
$10^{-6}$ on a case with known exact answer). Results:
\begin{itemize}
\item \textbf{No lattice correction exists for the model actually evaluated.} For the
isotropic/1D-channel continuum model underlying VF.2--VF.4, the exact Fermi-contour
integral equals the linearized value \emph{identically}; $C=50/(3\pi)$ is exact for that
model, and the quoted $5.339$ is not reproduced by Lorentzian-broadening, energy-bin, or
coarse-grid artifact classes (closest: $+0.76\%$ vs the required $+0.64\%$). The value is
struck from VF.4; the $+28\%$ threat to $M^*$ from that number is \textbf{retired}.
\item \textbf{New finding (van Hove obstruction).} On the honest $C_6$ triangular band the
$5/6$ edge-locking lands \emph{above} the van Hove energy in both natural directions
($E_F = +0.493, +0.515$ vs $E_{\rm vH} = +0.400$): the Fermi surface is disconnected
$K$-pockets, not a $\Gamma$-centered circle, and the isotropic linearization fails
qualitatively there ($-29\%$ / $+75\%$). $C = 50/(3\pi)$ is therefore a property of the
idealized isotropic model \emph{only}; a first-principles 2D realization in which the
$\Gamma$-circle picture survives at $5/6$ locking is now an explicit open requirement.
\item \textbf{New finding (unit mixing).} The formula $C = g\,(L_F/4\pi^2)(2/v_F)$
evaluates $L_F$ with $k_F$ as a bare fraction ($5/6$) but $v_F$ as a radian derivative (at
$k=5\pi/6$). The two unit-consistent evaluations give $C = 50/3$ or $50/(3\pi^2)$, which
miss $m_\tau$ by more than ten orders of magnitude; \emph{only} the mixed evaluation lands
on $m_\tau$. The unit convention is thus itself a load-bearing, currently unjustified
choice.
\end{itemize}
GAP-N4c status after the computation: the anchor survives the $5.339$ threat, but the
justification debt grows --- what would close the gap is (i) a first-principles 2D model
compatible with $\Gamma$-circle locking, (ii) a principled fixing of the momentum-unit
convention, and (iii) fixing of the BCS prefactor, the cutoff identification, and $g=4$.
